\documentclass[11pt]{article}
\usepackage[margin=1in]{geometry}
\usepackage{amsmath,amssymb,amsthm}
\usepackage[hidelinks]{hyperref}
\newtheorem{theorem}{Theorem}
\newtheorem{lemma}{Lemma}
\newcommand{\Res}{\operatorname*{Res}}
\title{Explicit polynomial solutions of Cohen's\\two difference equations}
\author{Principia Math In Progress}
\date{18 September 2026}
\begin{document}
\maketitle
\begin{abstract}
We prove the existence assertion in Cohen's Conjecture 4.2 for every
nonnegative integer parameter. An explicit finite polynomial is
constructed by a residue of a terminating hypergeometric expression.
Together with the previously observed degree, uniqueness, and reflection
properties, this settles both equations in MathDB 358171.
The proof uses two symbolic coefficient identities and does not depend
on a finite computational search.
\end{abstract}

\section*{Task 1: the conjecture and its complete proof}
For \(u\in\mathbb Z_{\ge0}\), put
\[
 N=u(u+3),\qquad n=N/2,\qquad
 \lambda=2u^2+8u+7,\qquad \mu=(2u+3)(N+1).
\]
The number \(N\) is even. Cohen's Conjecture 4.2 \cite{C},
on pages 7--8, asks for nonzero polynomial solutions of
\begin{equation}\label{first}
\begin{split}
 &(x+2u+4)(x+u+1)P(x+1)-(x-1)(x+u-1)P(x-1)\\
 &\hspace{25mm}=(\lambda x+\mu)P(x),
\end{split}
\end{equation}
\begin{equation}\label{second}
\begin{split}
 &(x+2u+4)(x+u+4)Q(x+1)-(x-1)(x+u+2)Q(x-1)\\
 &\hspace{25mm}=\bigl(\lambda x+(2u+3)(u+2)(u+3)\bigr)Q(x).
\end{split}
\end{equation}
The conjecture predicts degree \(N\) and uniqueness up to nonzero
scalar multiples for each equation. The degree, uniqueness conditional
on existence, and equivalence of the two equations are already due to
Peter Mueller \cite{M}; they are also credited on page 8 of \cite{C}.
We supply the missing existence proof and repeat those short arguments
to make the result self-contained.

\begin{theorem}\label{main}
For every nonnegative integer \(u\), the solution space in
\(\mathbb Q[x]\) of each equation \eqref{first} and \eqref{second}
is one-dimensional, and every nonzero solution has degree \(N=u(u+3)\).
\end{theorem}

\subsection*{The explicit construction}
Define
\[
 D=N+2u+5,\qquad K=N+u+3=(u+1)(u+3),\qquad
 e_k=\frac{(-n)_k(1/2)_k}{(-D/2)_k\,k!}\quad(0\le k\le n)
\]
and
\begin{equation}\label{H}
 H(z)=\sum_{k=0}^{n}e_k
       \left(z^{2k}+\frac{2k-K}{2k-D}z^{2k+1}\right)
      =\sum_{j=0}^{N+1}h_jz^j.
\end{equation}
All coefficients are rational and all denominators are nonzero:
\(D>N\), and \(-D/2\) is a negative half-integer.
The required polynomial is
\begin{equation}\label{res}
 P(x)=\Res_{t=0}
       e^{(x+u)t}H\bigl(\coth(t/2)\bigr)\,dt .
\end{equation}
This means the coefficient of \(t^{-1}\) in the indicated Laurent
series. Equivalently, it is the following completely finite expression:
\begin{equation}\label{finite}
 P(x)=\sum_{j=1}^{N+1}h_j
           \sum_{\ell=0}^{j-1}
           2^{\ell+1}\binom{j}{\ell+1}
                         \binom{x+u-1}{\ell},
 \qquad
 Q(x)=P(-x-2u-3).
\end{equation}
Here the generalized binomial is a polynomial in its upper argument.
The equivalence follows from \(w=e^t\), \(v=w-1\):
\[
 P(x)=\Res_{v=0}(1+v)^{x+u-1}
                 H\bigl((v+2)/v\bigr)\,dv.
\]
Expanding each finite power \((v+2)^j\) gives \eqref{finite}.

\subsection*{A polynomial differential equation}
\begin{lemma}\label{odeH}
The polynomial \eqref{H} satisfies
\begin{equation}\label{Heqn}
 z(z^2-1)H''(z)+(-Nz^2+3z+D)H'(z)-KH(z)=0.
\end{equation}
\end{lemma}
\begin{proof}
For any coefficients \(h_j\), the coefficient of \(z^j\) on the
left of \eqref{Heqn} is
\begin{equation}\label{rec}
 (j-1)(j-2-N)h_{j-1}+(3j-K)h_j
                       +(j+1)(D-j)h_{j+1}.
\end{equation}
Take coefficients outside \(0\le j\le N+1\) to be zero.
Write \(o_k=h_{2k+1}\), so that
\[
 h_{2k}=e_k,\qquad o_k=\frac{2k-K}{2k-D}e_k,\qquad
 \frac{e_{k+1}}{e_k}
 =\frac{(2k-N)(2k+1)}{2(k+1)(2k-D)}.
\]
The last relation also gives \(e_{n+1}=0\).
For \(j=2k\), substitution in \eqref{rec} reduces to
\[
 \bigl[2k(2k-2-K)+(6k-K)-(2k+1)(2k-K)\bigr]e_k=0.
\]
At \(k=0\) the first term is absent and the same identity applies.
For \(j=2k+1\), multiply \eqref{rec} by \((2k-D)/e_k\).
The result is
\begin{align*}
 &2k(2k-1-N)(2k-D)+(6k+3-K)(2k-K)\\
 &\hspace{35mm}+(D-2k-1)(2k-N)(2k+1).
\end{align*}
Using \(D=N+2u+5\) and \(K=N+u+3\), this expression is
\[
 u^2+3u-N=0.
\]
This includes \(k=n\), where the last term vanishes.
The only remaining possible coefficient is \(j=N+2\);
its factor \(j-2-N\) is zero. Higher coefficients vanish.
This proves the differential equation for every \(u\).
\end{proof}

\subsection*{Transporting the differential equation to a residue}
Set
\[
 g(w)=w^u H\!\left(\frac{w+1}{w-1}\right),
 \qquad f(t)=g(e^t)=e^{ut}H\bigl(\coth(t/2)\bigr).
\]
The chain rule transforms \eqref{Heqn} into
\begin{equation}\label{gode}
 w(w^2-1)g''+
 \bigl(-(3u+2)w^2+\lambda w+u-1\bigr)g'
 +\bigl(u(2u+3)w-\mu\bigr)g=0.
\end{equation}
Here are explicit coefficient checks for that transformation.
With \(z=(w+1)/(w-1)\), let \(a=w(w^2-1)\),
\(b=-(3u+2)w^2+\lambda w+u-1\), and
\(c=u(2u+3)w-\mu\).
After division by \(w^u\), the coefficients of \(H'',H',H\) are
\[
 a(z')^2,\qquad a\left(z''+\frac{2u}{w}z'\right)+bz',
 \qquad \frac{u(u-1)}{w^2}a+\frac{u}{w}b+c.
\]
Substitution of \(w=(z+1)/(z-1)\) gives, respectively,
\(z(z^2-1)\), \(-Nz^2+3z+D\), and \(-K\).
In turn, \(w=e^t\) converts \eqref{gode} into
\begin{equation}\label{fode}
 (e^t-e^{-t})f''+
 \bigl(\lambda-3(u+1)e^t+u e^{-t}\bigr)f'
 +\bigl(u(2u+3)e^t-\mu\bigr)f=0.
\end{equation}

We now verify \eqref{first} directly. Set
\[
 A(x)=x^2+(3u+5)x+(2u+4)(u+1),\qquad
 B(x)=x^2+(u-2)x+1-u.
\]
For any meromorphic germ \(F(t)\),
\[
 \Res_{t=0}(e^{xt}F(t))'\,dt=0,\qquad
 x\Res e^{xt}F\,dt=-\Res e^{xt}F'\,dt.
\]
Applying this identity twice shows that the left side of
\eqref{first} minus its right side equals the residue of
\[
 e^{xt}\left\{
 A(-\partial_t)(e^tf)-B(-\partial_t)(e^{-t}f)
                    +\lambda f'-\mu f\right\}dt .
\]
Expanding the braces gives exactly the left side of \eqref{fode}.
That expression vanishes identically, proving \eqref{first}.
All operations concern germs with finite-order poles, so no limiting
integral, summability, or exchange of infinite tails is needed.

\subsection*{Nonvanishing, degree, uniqueness, and the second equation}
The leading coefficient of \(H\) is
\[
 h_{N+1}=e_n\frac{u+3}{2u+5}\ne0.
\]
Since \(\coth(t/2)=2/t+O(t)\), equation \eqref{res} is a polynomial
of degree exactly \(N\), with leading coefficient
\[
 [x^N]P(x)=\frac{2^{N+1}h_{N+1}}{N!}\ne0.
\]
Its coefficients are rational, also directly evident from
\eqref{finite}. This proves the previously unresolved existence.

For clarity, the short uniqueness argument from \cite{M} is as follows.
If a nonzero polynomial of degree \(d\), with leading coefficient \(a\),
satisfies \eqref{first}, the coefficient of \(x^{d+1}\) in the residual is
\[
 a(2d+2u+7-\lambda)=2a(d-N).
\]
Thus \(d=N\). Two monic solutions would have a difference of smaller
degree, hence must coincide.
Finally, substitute \(y=-x-2u-3\) into \eqref{first} and multiply by
\(-1\). The two shifted coefficients interchange and
\[
 \lambda(2u+3)-\mu=(2u+3)(u+2)(u+3).
\]
Therefore \(Q(x)=P(-x-2u-3)\) satisfies \eqref{second}.
This substitution is invertible, so its solution space is also
one-dimensional. Theorem \ref{main} follows.

For example, the monic solutions of \eqref{first} at \(u=0,1\) are
\[
 1,\qquad x^4+\frac{15}{2}x^3+\frac{85}{2}x^2+105x+144.
\]
The construction also works for real \(u\ge0\) whenever \(u(u+3)\)
is an even integer; then the coefficients lie in \(\mathbb Q(u)\).
This follows from the same proof and earns no additional problem credit.

\small
\section*{Task 2: remove the logarithmic term in Cohen's application}
Using the explicit \(P\), settle the separate assertion on page 8
of \cite{C} that the associated alternating rational series has no
\(\log 2\) term. Verify the coprimality conditions needed for
Proposition 4.1 and obtain the initial constants in the resulting
continued fractions for \(\pi^2\).

\section*{Task 3: determine the arithmetic of the monic solution}
Determine the least integer clearing the denominators of the monic
\(P\), including its prime valuations as functions of \(u\).
Test and, if true, prove positivity of its ordinary coefficients and
irreducibility over \(\mathbb Q\). These properties are not assumed
in the existence proof.

\section*{Task 4: classify the full parameter locus}
For even \(N\ge0\), classify all pairs \(D,K\) for which
\[
 z(z^2-1)H''+(-Nz^2+3z+D)H'-KH=0
\]
has a polynomial solution of degree \(N+1\).
The construction identifies the component
\(D=N+2u+5,\ K=N+u+3,\ N=u(u+3)\).
Describe additional components, exceptional denominators, and their
corresponding difference equations.

\section*{Task 5: use the hypergeometric form to study zeros}
The even part of \(H\) is
\({}_2F_1(-N/2,1/2;-D/2;z^2)\).
Use this terminating representation and its contiguous odd part
to locate the zeros of \(H\) and \(P\), and to derive their scaled
zero distributions as \(u\) tends to infinity.
Determine whether these formulas yield faster exact construction
than solving for the coefficients by descending elimination.

\begin{thebibliography}{9}
\small
\bibitem{C}
H.~Cohen, \emph{Parametric Continued Fractions for \(\pi^2\),
\(\zeta(3)\), and other Constants}, arXiv:2304.11727v1 (2023),
\url{https://arxiv.org/abs/2304.11727}.
Conjecture 4.2 and the credited partial results are on pp.~7--8.
\bibitem{M}
P.~Mueller, answer to H.~Cohen,
\emph{Polynomial solutions to a difference equation},
MathOverflow, 10 April 2023,
\url{https://mathoverflow.net/questions/444367/}.
\end{thebibliography}
\end{document}
